% !TeX root = ../main.tex
\section{Additional Derivations}

\subsection{On Identifiability of the Zero-Drift Equilibrium}
\label{sec:theory_identifiability}

% \paragraph{Problem statement.}
In Sec.~\ref{sec:methodology}, we showed that anti-symmetry implies $p=q \Rightarrow \V(\x)\equiv \mathbf{0}$.
Here we investigate the converse: under what conditions does $\V(\x)\approx \mathbf{0}$ imply $p\approx q$?
Generally, this is not guaranteed for arbitrary vector fields.
However, we argue that for our specific construction, the zero-drift condition imposes strong constraints on the distributions.

% \paragraph{Full-support assumption.}
To avoid boundary issues, we assume that $p$ and $q$ have full support on $\mathbb{R}^d$ (e.g., via infinitesimal Gaussian smoothing). 
Consequently, ensuring the equilibrium condition $\V(\x) \approx \mathbf{0}$ for generated samples $\x \sim q$ effectively enforces $\V(\x) \approx \mathbf{0}$ for all $\x \in \mathbb{R}^d$.

\paragraph{Setup.}
Consider a general interaction kernel $K(\mathbf{x},\mathbf{y}^+,\mathbf{y}^-) \in \mathbb{R}^d$ and the drifting field
\begin{equation}
\mathbf{V}_{p,q}(\mathbf{x}) := \mathbb{E}_{\mathbf{y}^+\sim p,\ \mathbf{y}^-\sim q} \big[ K(\mathbf{x},\mathbf{y}^+,\mathbf{y}^-) \big].
\label{eq:V_general_K}
\end{equation}
We assume that $p$ and $q$ belong to a finite-dimensional model class spanned by a linearly independent basis $\{\varphi_i\}_{i=1}^m$:
\begin{equation}
p(\mathbf{y})=\sum_{i=1}^m a_i\,\varphi_i(\mathbf{y}),\qquad
q(\mathbf{y})=\sum_{i=1}^m b_i\,\varphi_i(\mathbf{y}),
\label{eq:pq_basis}
\end{equation}
where $\mathbf{a},\mathbf{b}\in\mathbb{R}^m$ are coefficient vectors.

\paragraph{Bilinear expansion over test locations.}
Consider a set of test locations (probes) $\mathcal{X} = \{\mathbf{x}_k\}_{k=1}^N$ with sufficiently large $N$ (e.g., $N \gg m^2$).
For each pair of basis indices $(i,j)$, we define the \emph{induced interaction vector} $\mathbf{U}_{ij} \in \mathbb{R}^{d{\times}N}$ by computing its column:
\begin{equation}
\mathbf{U}_{ij}[:, \mathbf{x}] \triangleq \iint K(\mathbf{x},\mathbf{y}^+,\mathbf{y}^-)\, \varphi_i(\mathbf{y}^+)\,\varphi_j(\mathbf{y}^-)\,d\mathbf{y}^+d\mathbf{y}^-
\end{equation}
evaluated at all $\mathbf{x} \in \mathcal{X}$.
Substituting the basis expansion into Eq.~\eqref{eq:V_general_K}, the drifting field evaluated on $\mathcal{X}$ (stored as a matrix $\mathbf{V}_{\mathcal{X}}$) is a bilinear combination:
\begin{equation}
\mathbf{V}_{\mathcal{X}} \triangleq \sum_{i=1}^m \sum_{j=1}^m a_i b_j \mathbf{U}_{ij}.
\end{equation}
Here, $\mathbf{V}_{\mathcal{X}} \in \mathbb{R}^{d{\times}N}$. At the equilibrium, we have $\mathbf{V}_{\mathcal{X}}=\textbf{0}$, which yields $dN$ linear equations.

\paragraph{Linear independence assumption.}
Our anti-symmetry condition implies that switching $p$ and $q$ negates the field. In terms of basis interactions, this means $\mathbf{U}_{ij} = -\mathbf{U}_{ji}$ (and consequently $\mathbf{U}_{ii} = \mathbf{0}$).
We make the \emph{generic non-degeneracy assumption}:
% \begin{center}
\textit{The set of vectors $\{\mathbf{U}_{ij}\}_{1 \le i < j \le m}$ is linearly independent in $\mathbb{R}^{dN}$.}
% \end{center}
This assumption requires the probes $\mathcal{X}$ and kernel $K$ to be non-degenerate; if all $\mathbf{x}$ yield identical constraints, independence would fail. For generic choices of $K$ and sufficiently diverse probes $\mathcal{X}$ with $dN\gg m^2$, such linear independence is a natural non-degeneracy condition.


\paragraph{Uniqueness of the equilibrium.}
The zero-drift condition $\mathbf{V}(\mathbf{x}) \equiv \mathbf{0}$ implies $\mathbf{V}_{\mathcal{X}} = \mathbf{0}$. Grouping terms by the independent basis vectors $\{\mathbf{U}_{ij}\}_{i<j}$, we have:
\begin{equation}
\sum_{1 \le i < j \le m} (a_i b_j - a_j b_i) \mathbf{U}_{ij} = \mathbf{0}.
\end{equation}
By the linear independence assumption, the coefficients must vanish: $a_i b_j - a_j b_i = 0$ for all $i,j$.
This implies that the vector $\mathbf{a}$ is parallel to $\mathbf{b}$ (i.e., $\mathbf{a} \propto \mathbf{b}$).
Since $p$ and $q$ are probability densities (implying $\int p = \int q = 1$), we must have $\mathbf{a} = \mathbf{b}$, and thus $p=q$.

\paragraph{Connection to the mean shift field.}
The mean-shift field fits this framework.
The update vector (before normalization) is
$\mathbb{E}_{p,q}[ k(\mathbf{x},\mathbf{y}^+) k(\mathbf{x},\mathbf{y}^-) (\mathbf{y}^+ - \mathbf{y}^-) ]$.
Assuming the normalization factors $Z_p$ and $Z_q$ are finite, the condition $\mathbf{V}(\mathbf{x}) = \mathbf{0}$ implies the numerator integral vanishes, which corresponds to an interaction kernel of the form:
\begin{equation}
K(\mathbf{x}, \mathbf{y}^+, \mathbf{y}^-) = k(\mathbf{x},\mathbf{y}^+)\, k(\mathbf{x},\mathbf{y}^-)\, (\mathbf{y}^+ - \mathbf{y}^-).
\end{equation}
This kernel generates the bilinear structure analyzed above.
Since we can choose $N$ such that $dN \gg m^2$, the dimension of the test space is much larger than the number of basis pairs. Thus, the linear independence of $\{\mathbf{U}_{ij}\}$ is expected to hold for generic configurations.
Finally, for general distributions $p$ and $q$, we can approximate them using a sufficiently large basis expansion, turning into $\tilde p$ and $\tilde q$. When the basis approximation is sufficiently accurate, $\tilde p \approx p$ and $\tilde q \approx q$, and the drift field $\mathbf{V}_{\tilde p, \tilde q}\approx \mathbf{V}_{p,q}\approx 0$. By the argument above, $\tilde p\approx \tilde q$, and thus $p\approx q$.


The argument above works for general form of drifting field, under mild anti-degeneracy assumptions.


\subsection{The Drifting Field of MMD}
\label{sec:theory_mmd_relation}

In principle, if a method minimizes a discrepancy between two distributions $p$ and $q$ and reaches minimum at $p=q$, then from the perspective of our framework, a drifting field $\V$ exists that governs sample movement: we can let $\V{\propto}{-}\frac{\partial\mathcal L }{\partial {\x}}$, which is zero when $p=q$.
We discuss the formulation of this $\V$ for a loss based on Maximum Mean Discrepancy (MMD) \cite{li2015generative,dziugaite2015training}.

\newcommand{\kkk}{\tilde{k}}
\paragraph{Gradients of Drifting Loss.}
With $\x=f_\theta(\e)$, our drifting loss in Eq.~(\ref{eq:train_loss}) can be written as:
\begin{equation}
\mathcal{L}
=
\mathbb{E}_{\x{\sim}q}
[
\mathcal{L}(\x)
]
=
\mathbb{E}_{\x{\sim}q}
\Big[
\big\|
\x
-
\text{\texttt{sg}}\big(
\x
+
\V(\x)
\big)
\big\|^2
\Big],
\end{equation}
where ``\texttt{sg}'' is short for stop-gradient. The gradient w.r.t. the parameters $\theta$ is computed by:
\begin{equation}
\frac{\partial\mathcal L }{\partial {\theta}} =
\mathbb{E}_{\x{\sim}q}\Big[
\frac{\partial\mathcal L(\x) }{\partial {\x}} \frac{\partial\x}{\partial\theta}
\Big].
\end{equation}
where
$
\frac{\partial\mathcal L(\x) }{\partial {\x}}{=}2
(
\x
{-}
\text{\texttt{sg}}(
\x
{+}
\V(\x)
)
)
{=}{-}2\V(\x)
$. This gives:
\begin{equation}
\label{eq:loss_2V}
\V(\x) = -\frac{1}{2}\frac{\partial\mathcal L(\x) }{\partial {\x}}
\end{equation}
We note that this formulation is general and imposes no constraints on $\V$, except that $\V=\textbf{0}$ when $p=q$.

Our method does not require $\mathcal{L}$ to define a discrepancy between $p$ and $q$. However, for other methods that depend on minimizing a discrepancy $\mathcal{L}$, we can induce a drifting field via (\ref{eq:loss_2V}). This is valid if $\mathcal{L}$ is minimized when $p=q$.

\paragraph{Gradients of MMD Loss.}
\newcommand{\kk}{\xi}
\newcommand{\LMMD}{\mathcal{L}_{\text{MMD}^2}}
In MMD-based methods (\eg, \citealt{li2015generative}), the difference between two distributions $p$ and $q$ is measured by squared MMD:
\begin{equation}
\label{eq:mmd2_def}
\begin{aligned}
\LMMD(p,q)
= &
\mathbb{E}_{\x,\x'\sim q}[\kk(\x,\x')]
-2\,\mathbb{E}_{\y\sim p,\;\x\sim q}[\kk(\y,\x)]
\\
+ &~const.
\end{aligned}
\end{equation}
Here, the constant term is $\mathbb{E}_{\y,\y'\sim p}[\kk(\y,\y')]$, which depends only on the target distribution $p$ and remains unchanged. $\kk$ is a kernel function. 

Consider $\x=f_\theta(\e)$ with $\e\sim p_\e$. The gradient estimation performed in \cite{li2015generative} corresponds to:
\begin{equation}
\frac{\partial\LMMD}{\partial {\theta}} =
\mathbb{E}_{\x{\sim}q}\Big[\frac{\partial\LMMD(\x)}{\partial {\x}} \frac{\partial\x}{\partial\theta}
\Big]
\end{equation}
where the gradient w.r.t $\x$ is computed by:
\begin{equation}
\frac{\partial\LMMD(\x)}{\partial {\x}}
=
2\mathbb{E}_{\x'{\sim}q} \Big[\frac{\partial \kk(\x,\x')}{\partial{\x}}\Big]
-
2\mathbb{E}_{\y{\sim}p} \Big[\frac{\partial \kk(\x,\y)}{\partial{\x}}\Big].
\end{equation}
Using our notation of positives and negatives, we rename the variables and rewrite as:
\begin{equation}
\label{eq:mmd_eq4}
\frac{\partial\mathcal L_{\text{MMD}^2}(\x)}{\partial {\x}}
=
2\mathbb{E}_{\yn{\sim}q} \Big[\frac{\partial \kk(\x,\yn)}{\partial{\x}}\Big]
-
2\mathbb{E}_{\yp{\sim}p} \Big[\frac{\partial \kk(\x,\yp)}{\partial{\x}}\Big].
\end{equation}
Comparing with Eq.~(\ref{eq:loss_2V}), we obtain:
\begin{equation}
\label{eq:mmd_Vmmd}
\V_\text{MMD}(\x)
\triangleq
\mathbb{E}_{\yp{\sim}p} \Big[\frac{\partial \kk(\x,\yp)}{\partial{\x}}\Big]
-
\mathbb{E}_{\yn{\sim}q} \Big[\frac{\partial \kk(\x,\yn)}{\partial{\x}}\Big]
\end{equation}
This is the underlying drifting field that corresponds to the MMD loss $\LMMD$.




For a radial kernel $\kk(\x, \y)=\kk(R)$ where  $R=\|\x - \y\|^2$, the gradient of kernel is:
\begin{equation}
\frac{\partial\kk(\x, \y)}{\partial\x} = 
2\kk'(\|\x - \y\|^2)(\x - \y)
\end{equation}
where $\kk'$ is the derivative of the function $\kk(R)$.
Accordingly, Eq.~(\ref{eq:mmd_Vmmd}) becomes:
\begin{equation}
\begin{aligned}
\label{eq:mmd_Vmmd2}
\V_\text{MMD}(\x)
= &
\mathbb{E}_{\yp{\sim}p} \Big[
2\kk'(\|\x - \yp\|^2)(\x - \yp)
\Big].
\\ - &
\mathbb{E}_{\yn{\sim}q} \Big[
2\kk'(\|\x - \yn\|^2)(\x - \yn)
\Big]
\end{aligned}
\end{equation}
In \cite{li2015generative}, the Gaussian kernel is used: $\kk(\x, \y) = \exp (-\frac{1}{2\sigma^2} \|\x{-}\y\|^2)$, leading to $\kk'(\|\x{-}\y\|^2)=-\frac{1}{2\sigma^2}\exp (-\frac{1}{2\sigma^2} \|\x{-}\y\|^2)$.

\paragraph{Relations and Differences.}
When using our definition of $\V = \V^{+} - \V^{-}$ (\ie, Eq.~(\ref{eq:V_pos_neg})), we have:
\begin{equation}
\begin{aligned}
\label{eq:mmd_eq6}
\V(\x)
= & 
\mathbb{E}_{\yp{\sim}p} \Big[
\tilde{k}(\x, \yp)(\yp - \x)
\Big] \\
- &
\mathbb{E}_{\yn{\sim}q} \Big[
\tilde{k}(\x, \yn)(\yn - \x)
\Big]
\end{aligned}
\end{equation}

Comparing (\ref{eq:mmd_Vmmd2}) with (\ref{eq:mmd_eq6}), we show that the underlying kernel used to build the drifting field of MMD is:
\begin{equation}
\label{eq:mmd_eq7}
\tilde{k}_\text{MMD}(\x, \y)=-2\kk'(\|\x - \y\|^2).
\end{equation}
When $\kk$ is a Gaussian function, we have: $\tilde{k}(\x, \y)=\frac{1}{\sigma^2}\exp (-\frac{1}{2\sigma^2} \|\x{-}\y\|^2)$.
Without normalization, the resulting drift no longer satisfies the assumptions underlying Alg.~\ref{alg:drift_code}, and the mean-shift interpretation breaks down.

As a comparison, our general formulation enables to use \textit{normalized} kernels:
\begin{equation}
\tilde{k}(\x, \y)=\frac{1}{Z(\x)}k(\x, \y)=\frac{1}{\mathbb{E}_\y[k(\x, \y)]}k(\x, \y),
\end{equation}
where the expectation is over $p$ or $q$.
Only when we use normalized kernels, we have (see Eq.~(\ref{eq:VZZ})):
\begin{equation}
\V(\x)
= {\mathbb{E}_{p,q}\!\left[\tilde{k}(\x,\yp)\tilde{k}(\x,\yn)(\yp - \yn)\right]},
\end{equation}
on which our Alg.~\ref{alg:drift_code} is based.

Given this relation, we summarize the key differences between our model and the MMD-based methods as follows:
\begin{enumerate}
[itemsep=1pt, topsep=1pt, parsep=1pt, partopsep=1pt, label=(\roman*)]
    \item Our method is formulated around the drifting field $\V$, which is more flexible and general.
    \item Our method supports and leverages \textit{normalized} kernels $\frac{1}{Z}k(\x, \y)$ that cannot be naturally derived from the MMD perspective.
    % \item Our method does not require the kernel to be the derivative of another function (\cf~(\ref{eq:mmd_eq7})). This requirement cannot be satisfied if $\tilde{k}$ is a normalized kernel.
    \item Our $\V$-centric formulation enables a flexible step size for drifting (\ie, $\x \leftarrow \x{+}\eta\V$) and therefore naturally supports $\V$-normalization (see \ref{sec:appendix_feature_and_loss_normalization}).
    \item Our $\V$-centric formulation allows the equilibrium concept to be naturally extended to support CFG, whereas a CFG variant for MMD remains unexplored.
\end{enumerate}
In summary, although a special case of our method reduces to MMD, our $\V$-centric framework is more general and enables unique possibilities that are important in practice.
In our experiments, we were not able to obtain reasonable results using the MMD framework.

